\documentclass[12pt]{article}
\usepackage[utf8]{inputenc}
\usepackage[T1]{fontenc}
\usepackage{hyperref}
\title{Falam Chin Language Resources for K-12 Curriculum Development}
\author{Compiled by Claude Sonnet 5.5 (claude-sonnet-5-5) for Andr\'as Kornai}
\date{}
\begin{document}
\maketitle

\section{Introduction}
The code flm is a retired ISO 639-3 code. The search results I found describe Falam Chin under the code cfm, so the resources below are for Falam Chin generally. I did not verify that the retired code flm and cfm cover exactly the same variety.

\section{Primary Educational Resources}
No primers or school textbooks were found in search results.

\section{Linguistic and Reference Materials}
\begin{itemize}
\item Falam Chin--English dictionary on Webonary (compiled by Than Bil Luai and Timothy Mang, 2018, with Bibles International and SIL; about 13,000 entries): \url{https://www.webonary.org/chinfalam/language/map/}
\item Endangered Languages Project entry for Falam Chin: \url{https://www.endangeredlanguages.com/node/39902}
\end{itemize}

\section{Religious and Community Resources}
\begin{itemize}
\item Falam Chin Bible, 2009 edition, Digital Bible Society: \url{https://dbs.org/bibles/FLMBIV}
\end{itemize}

\section{Digital and Audio Resources}
The Webonary dictionary above is the only digital resource found. No audio or video was found.

\section{Curriculum Development Context}
A dictionary and a complete Bible exist, which give a base for vocabulary and reading texts. Graded readers, textbooks and teacher materials still need to be developed. No school programme was found in search results.

\end{document}